Dans les r\'eseaux logistiques mutualis\'es, l'urgence qu'un op\'erateur d\'eclare diverge
souvent de celle que l'on peut calculer \`a partir de l'\'etat op\'erationnel. La
sur-d\'eclaration gaspille de la capacit\'e partag\'ee~; la sous-d\'eclaration dissimule un
risque de rupture. Aucun jumeau num\'erique de cha\^ine d'approvisionnement existant ne
mesure cette divergence ni ne la propage \`a l'\'echelle d'un r\'eseau de fournisseurs.

Ce m\'emoire en construit un et le met \`a l'\'epreuve. L'urgence r\'eelle est agr\'eg\'ee \`a
chaque n\oe{}ud \`a partir de six blocs de KPI par un OU probabiliste pond\'er\'e, puis
propag\'ee vers l'amont sur un graphe orient\'e acyclique~; l'urgence d\'eclar\'ee est
\'elicit\'ee par un questionnaire de comparaisons par paires soumis \`a un contr\^ole de
coh\'erence, puis propag\'ee vers l'aval~; leur \'ecart est not\'e sur $[0,100]$ par une
p\'enalit\'e asym\'etrique issue de la th\'eorie des perspectives. Le syst\`eme a \'et\'e
\'evalu\'e par deux campagnes pr\'e-enregistr\'ees dont les d\'eclarants sont des mod\`eles de
langage~: un rejeu en dix-neuf tours de la p\'enurie de semi-conducteurs de 2020 \`a 2022 sur
huit n\oe{}uds, avec un bras t\'emoin en aveugle et un bras assist\'e par pr\'evision sur un
sc\'enario gel\'e~; et la base d'approvisionnement d'un avionneur, quinze n\oe{}uds sans
aucun choc exog\`ene, dont la v\'erit\'e jalon d\'ecoule de la trajectoire des indicateurs au
lieu d'\^etre \'ecrite \`a la main.

Trois r\'esultats en d\'ecoulent. Comme instrument pour d\'ecider quel fournisseur regarder,
l'\'ecart fonctionne, et t\^ot~: sur les six tours pr\'ec\'edant l'irruption publique de la
crise, il ne signale que des n\oe{}uds qui rateront un engagement de livraison, trois des
quatre qui finiront par le faire, et le plus mal not\'e est signal\'e huit tours avant son
\'echec alors que son propre responsable d\'eclare une urgence de $0{,}001$. La seconde
cha\^ine borne cette affirmation au lieu de la confirmer~: l\`a o\`u les d\'eclarants jugent
correctement leur \'etat et diff\`erent simplement l'alerte, l'\'ecart ne d\'esigne rien.
L'instrument d\'etecte donc une m\'eprise, non une r\'etention. Comme instrument pour
attacher une probabilit\'e calibr\'ee \`a une semaine nomm\'ee, il \'echoue~: \emph{skill} de
Brier compris entre $-3{,}4$ et $-6{,}9$ face \`a un pr\'evisionniste annon\c{c}ant le taux de
base, et vingt-et-une pr\'evisions confiantes dont aucune ne s'est r\'ealis\'ee. Un seul bloc
l'explique~: le bloc temporel lit une distribution de d\'elai qu'aucun choc n'atteint jamais.
Le correctif n'est pas celui qu'on croit, puisque le temps perdu doit s'ajouter \`a la date
d'ach\`evement et non s'y multiplier, et il divise par deux le d\'eficit de Brier sans le
rendre positif.

Deux constats de m\'ethode d\'epassent ce syst\`eme. La d\'efinition de l'issue d\'eplace la
calibration rapport\'ee d'un facteur cinq sans toucher \`a la discrimination. Et afficher une
pr\'evision au d\'eclarant d\'eplace l'autorit\'e du terrain vers le mod\`ele, dans le sens
o\`u pointe le mod\`ele, sur une couche qu'un rapport calendaire \`a deux nombres surclasse
\`a six horizons sur huit. Ce qui borne chacun de ces taux est l'\'echantillon et non la
construction~: deux sc\'enarios, vingt-trois n\oe{}uds, sept positifs, trente-trois jalons
observables.

\vspace{0.4cm}

\noindent\textbf{Mots-cl\'es~:} jumeau num\'erique de cha\^ine d'approvisionnement,
ad\'equation d'urgence, biais cognitif, aide multicrit\`ere \`a la d\'ecision, calibration de
pr\'evisions, biais d'automatisation, jeu s\'erieux.
